% Section to follow paper/math_extension.tex. Requires amsmath, amsthm, amssymb.
% Numbers in the remarks come from experiments/aggregate.py; none are entered by hand.

\section{Stall-gated autoregulation}

\subsection{The layer as a matrix product}

Collect the per-head outputs $O_h \in \mathbb{R}^{d_k \times N}$ into the stacked matrix
$O = [\,O_1^{\top} \cdots O_H^{\top}\,]^{\top} \in \mathbb{R}^{H d_k \times N}$, and let
$G = \operatorname{diag}(g) \in \mathbb{R}^{H \times H}$ be the supply field. The gated
layer of Definition~\ref{def:supply} is then a single product,
\begin{equation}
  \label{eq:kron}
  \mathcal{T}_g(X) \;=\; W_O \,\bigl( G \otimes I_{d_k} \bigr)\, O,
  \qquad \operatorname{tr} G = H,
\end{equation}
with $\otimes$ the Kronecker product. Conservation is the trace constraint, and the
dense layer is $G = I_H$. Writing it this way makes the reduction transparent: for a
perfused set $S$ with $|S| = B$ and zero leak, $g = (H/B)\mathbf{1}_S$ gives
$G = (H/B) P_S$ with $P_S$ the diagonal $0/1$ projector onto $S$, so
\begin{equation}
  \mathcal{T}_g(X)
  \;=\; \frac{H}{B}\, W_O \,(P_S \otimes I_{d_k})\, O
  \;=\; \widetilde{W}_{O,S}\, O_S ,
\end{equation}
which is Proposition~\ref{prop:reduction}: the layer \emph{is} a $B$-head layer, and
\begin{equation}
  \operatorname{rank} \mathcal{T}_g(X) \;\le\; B \, d_k .
\end{equation}
The projector $P_S$ is what the supply rule chooses; everything below concerns how.

\subsection{The controller as a scalar dynamical system}

Let $B(\kappa) = |\{h : d_h > \kappa\}|$, nonincreasing in $\kappa$, and let
$L_\infty(B)$ be the loss the model converges to when held at $B$ perfused heads,
nonincreasing in $B$. Plain autoregulation (Definition~\ref{def:autoreg}) is the scalar
recursion
\begin{equation}
  \label{eq:plain}
  \kappa_{t+1} \;=\; \Pi\bigl(\kappa_t + \eta\,e_t\bigr),
  \qquad e_t \;=\; \log L^\star - \log \widehat{L}_t .
\end{equation}
The difficulty is that $e_t$ is built from the \emph{instantaneous} loss
$\widehat{L}_t$, whereas Theorem~\ref{thm:autoreg} is a statement about
$L_\infty(B)$. Early in training $\widehat{L}_t \gg L_\infty(B_t)$ for every $B_t$, so
$e_t < 0$ persistently and \eqref{eq:plain} integrates downward,
\begin{equation}
  \kappa_t \;\approx\; \kappa_0 - \eta \sum_{s<t} |e_s| \;\longrightarrow\; \kappa_-,
\end{equation}
driving $B \to H$ before the learner has had a chance to reveal what $B$ was
sufficient. This is not a tuning failure but a specification failure: the loop is
controlling the wrong quantity.

\subsection{The stall gate}

Let $\bar{L}_t$ be a slower exponential average of the loss than $\widehat{L}_t$, and
define the \emph{relative improvement rate}
\begin{equation}
  \rho_t \;=\; \frac{\bar{L}_t - \widehat{L}_t}{\bar{L}_t} \;\;\ge 0 \text{ while the
  learner is still descending.}
\end{equation}

\begin{definition}[Stall-gated autoregulation]
\label{def:stall}
Fix $\varepsilon > 0$. The loop is
\begin{equation}
  \label{eq:stall}
  \kappa_{t+1} \;=\; \Pi\bigl(\kappa_t + \eta\,\phi_t\, e_t\bigr),
  \qquad
  \phi_t \;=\;
  \begin{cases}
    0, & e_t < 0 \text{ and } \rho_t > \varepsilon,\\
    1, & \text{otherwise.}
  \end{cases}
\end{equation}
Dilation is admitted only when the deficit is \emph{chronic}: the loss is above target
\emph{and} has stopped improving. Constriction, $e_t \ge 0$, is never gated.
\end{definition}

The asymmetry is the content. In tissue the vasodilatory signal is the accumulation of a
metabolic byproduct, which requires the shortfall to persist; a transient dip does not
dilate a vessel. Equation \eqref{eq:stall} is that, and its effect is to change the
controlled variable.

\begin{proposition}[The gate restores the theorem's hypothesis]
\label{prop:stall}
Suppose the learner, held at a fixed perfused set, satisfies
$\widehat{L}_t \downarrow L_\infty(B)$ with $\rho_t \to 0$. Then along the trajectory of
\eqref{eq:stall} every dilation step occurs at a time $t$ with $\rho_t \le
\varepsilon$, so the error acting on $\kappa$ is
\begin{equation}
  e_t \;=\; \log L^\star - \log \widehat{L}_t
  \;=\; \log L^\star - \log L_\infty(B_t) \;+\; O(\varepsilon),
\end{equation}
and \eqref{eq:stall} is, to first order in $\varepsilon$, the recursion
\eqref{eq:plain} driven by $L_\infty(B_t)$ rather than by $\widehat{L}_t$. Its fixed
points are therefore those of Theorem~\ref{thm:autoreg}, and the selected perfusion is
$B^\star(L^\star) = \min\{B : L_\infty(B) \le L^\star\}$.
\end{proposition}

\begin{proof}
By construction $\phi_t = 0$ whenever $e_t < 0$ and $\rho_t > \varepsilon$, so any step
that decreases $\kappa$ has $\rho_t \le \varepsilon$. Since $\rho_t$ measures the
relative descent of $\widehat{L}$ toward its limit, $\rho_t \le \varepsilon$ gives
$\widehat{L}_t = L_\infty(B_t)(1 + O(\varepsilon))$, and the logarithm converts the
multiplicative error into the additive one stated.
\end{proof}

\begin{remark}[Two timescales, enforced rather than hoped for]
\label{rem:timescales}
Proposition~\ref{prop:stall} is a separation-of-timescales argument. The plain loop
attains the same separation only by taking $\eta$ small enough that the controller is
slower than the learner everywhere, which is a single global constant asked to respect a
rate that varies over training. The stall gate enforces the separation \emph{adaptively}:
the loop waits exactly as long as the learner needs at that point in training, and no
longer. The control that distinguishes the two is plain autoregulation at
$\eta/10$; if sluggishness alone sufficed, that arm would match the gated one.
\end{remark}

\begin{remark}[What must be validated]
\label{rem:validate}
The gate introduces $\varepsilon$, so the obvious objection is that the tuned target has
merely been replaced by a tuned threshold. Proposition~\ref{prop:stall} predicts
otherwise: $\varepsilon$ enters only as an $O(\varepsilon)$ perturbation of the control
error, so the selected $B$ should be insensitive to it over any range in which the
learner's descent is resolvable, whereas $L^\star$ enters the fixed point directly.
Three experiments follow, and all three are necessary:
\begin{enumerate}
  \item $B$ against $k^\star$ over $k^\star \in \{2,3,4,6,8\}$, the tracking test;
  \item $B$ against $L^\star$ over a $40\times$ range, the test that plain
        autoregulation failed, where the prediction is now that the staircase gap of
        Theorem~\ref{thm:autoreg} is recovered;
  \item $B$ against $\varepsilon$, where the prediction is insensitivity, and where a
        strong dependence would show the hyperparameter has only moved.
\end{enumerate}
A fourth, the second benchmark \texttt{qsa\_cross}, tests whether the result survives a
change in \emph{why} the circuit has a minimum size: there $k^\star$ is set by capacity,
$k^\star = \lceil q \log_2 N / d_k \rceil$, rather than by function.
\end{remark}
